Parameter-efficient fine-tuning (PEFT) adapts pretrained models with small trainable modules, but most adapters do not explicitly control how updates interact with the singular structure of the frozen weights. We study PEFT in the pretrained singular coordinate system and introduce a spectrum-aware adaptation framework that separates preserved pretrained components from task-specific adaptation. The framework first exposes a fixed-basis bottleneck: when adaptation is restricted to a prescribed low-spectrum basis, performance depends on how well that basis aligns with the desired task-specific perturbation. To address this limitation, we introduce compact orthogonal adaptation within the selected low-spectrum subspace, allowing directional flexibility while keeping the parameter cost tied to the spectral budget. We further instantiate these ideas with unit-invariant diagonal scalers, yielding practical adapters that provide coordinate-level flexibility while allowing us to bound scaler-induced mixing between preserved leading components and adapted tail components. Empirically, our proposed variants achieve the strongest average performance among the compared baselines on RoBERTa-base and RoBERTa-large GLUE, with the linear variant achieving strong results using substantially fewer trainable parameters. Additional generation results, retention-oriented experiments, and a diagnostic
soft-mixing ablation support the broader preservation-adaptation perspective.
